\section{Results \& Evaluation}
\label{sec:results_evaluation}

Having established a rigorous, multi-faceted benchmarking methodology in Section \ref{sec:experimental_setup},
this final section presents the empirical evaluation of the \texttt{ULTRA} generation of the \texttt{DetectorNX} architecture against its \gls{cnn} counterpart.
The following results were conducted on the pre-processed \texttt{PRECOMPUTED\_REPLICA\_10K\_V2}, using a validation split to ensure strict comparative integrity.

% Rather than relying on a singular metric of spatial accuracy, the following analysis systematically deconstructs the performance-efficiency frontier of the neuromoprhic paradigm.
% By sequentially examining 
% global detection parity, dynamic energy consumption, quantization-induced scaling errors, and predictive calibration, 
% this evaluation aims to quantify the precise trade-offs inherent in transitioning from 32-bit floating-point convolutions to 1-bit, 
% event-driven spiking dynamics for high-definition automotive perception.